\section{Identification of the complete cluster collections}
\label{sec:identification}

The stopped carpets determine the odd holes and their relative height
jumps. They do not determine which holes belong to one current cluster.
We identify that partition first. Finite families of visible holes then
transfer the lattice count bounds to all generations of the canonical
construction. The remaining task is to recover every point of a cluster,
including points that are not detected by an odd hole.

\subsection{Marked odd holes and their supplying clusters}

Fix a countable dense set \(\mathcal Q\subset D\). For each \(q\in\mathcal Q\)
choose a nearest lattice representative on the relevant \(W\) lattice,
using one deterministic rule to break ties. We record odd holes by a
family label \(\iota\in\{\mathrm f,\mathrm w\}\), a mark \(q\), and an
odd depth \(r\geq1\). The label \(\mathrm f\) refers to the original free
dual current; \(\mathrm w\) refers only to the free primal descendants
after the wired boundary cluster has been exposed.

At depth zero, \(V_{\delta,0}^{\mathrm f}(q)\) is the original free
chamber, with wall value \(b_{\delta,0}^{\mathrm f}(q)=0\).
For \(\iota=\mathrm w\), it is the component face left by the wired
layer that contains the representative of \(q\), with its observed
value \(b_{\delta,0}^{\mathrm w}(q)\in\{-a,a\}\).
Given a valid \(V_{\delta,r-1}^{\iota}(q)\), expose its first odd
carpet and let \(V_{\delta,r}^{\iota}(q)\) be the unsearched hole
containing the representative, with observed wall value
\(b_{\delta,r}^{\iota}(q)\). Write \(C_{\delta,r}^{\iota}(q)\)
for the current cluster supplying that odd rim, and put
\[
 \epsilon_{\delta,r}^{\iota}(q)
   =\frac{b_{\delta,r}^{\iota}(q)
               -b_{\delta,r-1}^{\iota}(q)}{2a}\in\{-1,1\}.
\]
If the representative is outside the chamber, lies on the revealed
carpet, or no such hole remains, give this observation and all its
descendants the cemetery value \(\dagger\), and set its sign to zero.
Validity, odd depth, and the supplying cluster depend only on the
unsigned data \((N,s)\), not on the free transverse colors \(d\).

The continuum observations \(V_r^\iota(q)\), \(b_r^\iota(q)\), and
\(\epsilon_r^\iota(q)\) are defined in the same way from the canonical
odd carpets of \cref{lem:canonical-geometry}. At depth zero the free
parent is \(D\), while the wired parent is the component of
\(D\setminus C_\partial\) containing \(q\), with its label \(\pm a\).
Use the same cemetery convention if a mark lies on a continuum carpet.
For a valid observation, \(C_r^\iota(q)\) denotes the retained canonical
cluster whose odd hole is \(V_r^\iota(q)\). The one-arm estimates in
the proof of \cref{lem:cut-neighborhoods} make every fixed mark miss
the limiting carpets almost surely. Every canonical odd hole contains a mark, so
these countably many observations see all odd holes at every finite
depth.

Here and below convergence of a marked hole is \emph{pointed}:
the closed carpet and its complement determine the limiting component
containing the mark, and every compact subset of that component is
eventually inside the discrete component. We do not assert Hausdorff
convergence of the closures of individual holes. In particular, one
discrete component may split into several pointed limiting components.

Abbreviate an address \((\iota,q,r)\) by \(i\), and write \(V_\delta(i)\),
\(C_\delta(i)\), and \(\epsilon_\delta(i)\), with the analogous notation
in the limit. For two addresses in the same family, define
\(i\sim_\delta j\) when both are valid and
\(C_\delta(i)=C_\delta(j)\); otherwise the relation is false.
Define \(i\sim j\) by the same rule for the supplying canonical
clusters, with a false relation whenever either address is invalid.

Fix any mesh subsequence. The whole-domain bound
\cref{cor:global-cluster-counts} and \cref{lem:matching-topology}
give compact containment of the two encoded collection factors
\(F(\cC_\delta^{\rm f})\) and \(F(\cC_\delta^{\rm w})\) in the
locally finite part of \(\mathfrak X\), the compact hyperspace over
the fixed enclosing square \(X\).
The height is tight by \cref{thm:flat-gff}. We pass to a joint further
subsequence retaining the closed carpets, their marked pointed
components and labels, the marked supplier supports in
\(\mathscr H\cup\{\dagger\}\) with an isolated cemetery value, and
the bits \(i\sim_\delta j\) for every fixed pair of addresses.
These supplier coordinates are compact;
they are separate from the contributor lists introduced later.
Write \(F_*^\iota\) for the locally finite outer Hausdorff limits of
the encoded collections. The relation coordinates are finite-valued;
on a joint representation they are eventually constant, with limiting
relation denoted by \(\sim_*\). This retains the finite identity
relations needed below, not numerical cluster identifiers across meshes.
By \cref{prop:stopped-tvs} and the one-arm estimate, for every finite
family of addresses the validity indicators and the signs in
\(\{-1,0,1\}\) converge jointly to their canonical counterparts.
In particular all bounded products of these signs converge in
expectation.

The use of \(X\) is necessary because an admissible polygon may
protrude outside \(D\). For each compact set disjoint from \(\partial D\),
uniform convergence of the boundary parametrizations preserves its
winding number for all sufficiently small mesh. Hence the polygonal
closures eventually lie in every fixed neighborhood of \(\overline D\),
and every Hausdorff limit of lattice supports lies in \(\overline D\).

We will use one elementary property of pointed convergence. If a sealed
discrete hole has pointed limit \(V\), every point of \(\partial V\)
is approached by its sealing boundary. Otherwise a ball around such a
point would avoid the discrete boundaries along a subsequence. Since
it meets a compact subset of \(V\), that connected ball would then lie
in the discrete hole, enlarging its pointed limit across the boundary.
The sealing boundary lies within \(O(\delta)\) of its supplying current
cluster. Consequently any Hausdorff limit of those suppliers contains
\(\partial V\); if the supplier has been revealed at a stop, the limiting
explored set contains \(\partial V\) as well.

\begin{lemma}[The restricted field determines its flat label]
\label{lem:recover-flat-label}
Let \(O\) be a random bounded simply connected domain with a marked
point \(q\in O\), and let \(f\) be a distribution on \(O\).
Suppose \(O,q,m\) are measurable in a sigma-field \(\mathcal E\) and,
conditionally on \(\mathcal E\), one has \(f=m+h^O\), where \(m\in\R\)
and \(h^O\) is a zero-boundary GFF. There is a measurable function of
\((O,q,f)\) alone that equals \(m\) almost surely. Thus two such conditional descriptions
of the same pointed domain and restricted field have the same label.
\end{lemma}

\begin{proof}
Let \(\phi_O:\mathbb D\to O\) be normalized by
\(\phi_O(0)=q\) and \(\phi'_O(0)>0\). Rational polygonal inner
exhaustions, normal-family compactness, and uniqueness of the normalized
limit make \(\phi_O\) measurable in \((O,q)\).
Push uniform angular probability measure \(\dd\vartheta/(2\pi)\) on
\(|z|=r\) forward by \(\phi_O\). Conformal invariance of the Green
kernel gives covariance \(-\log\max(r,s)\) for the zero-boundary
averages at radii \(r,s\). The radius-\(r\) average of \(f\) thus has
mean \(m\) and variance \(-\log r\).

For \(r_j=1-2^{-j}\), approximate the angular measure by a smooth
radial probability test \(\nu_j\) compactly supported in the disk,
with squared Green-energy error at most \(2^{-j}\).
Such choices are deterministic: the covariance
\(-\log\max(r,s)\) is continuous near each positive \(r_j\), so
radial smoothing converges in energy. Push \(\nu_j\) forward by
\(\phi_O\) to obtain a smooth compactly supported test \(\psi_j(O,q)\).
This test is measurable in the pointed domain, and evaluation
\(f(\psi_j(O,q))\) is measurable in \((O,q,f)\). Its conditional
mean is \(m\), and its conditional variance is at most
\[
 2(-\log r_j)+2^{1-j}.
\]
The sum of these variances is finite. Chebyshev's inequality and
Borel--Cantelli therefore give \(f(\psi_j(O,q))\to m\) almost surely.
Define the function in the statement by this limit when it exists,
and by zero otherwise. The construction uses conformal maps only on
compact subsets of the disk, so no boundary extension of \(O\) is
needed. The conformal identities are recalled in
\cite[Chapter~6, Sections~1.1 and~5.2]{Ahlfors1979}.
\end{proof}

\begin{proposition}[The visible partition]
\label{prop:partition}
Almost surely, \(i\sim_*j\) if and only if \(i\sim j\), simultaneously
for every pair of marked addresses in either current family. Their
relative signs also agree. Each retained canonical cluster \(C\) has
a marked address \(i\) such that every Hausdorff limit of
\(C_\delta(i)\) contains \(C\). For every finite family of distinct
canonical clusters, their chosen supplying lattice clusters are
eventually distinct.
\end{proposition}

\begin{proof}
We first show \(i\sim j\Rightarrow i\sim_*j\). All odd holes of one
retained canonical cluster have the same odd depth \(r\) and the same
parent at depth \(r-1\): entering any odd hole increases the depth, and
later clusters lie inside it. Let \(V_1,V_2\) be two marked odd holes
of such a cluster \(C\), and suppose that their lattice suppliers are
distinct along the retained subsequence.
If \(V_1=V_2\), a compact path inside that hole already puts both
marks in the same discrete hole for all sufficiently small mesh.
We may therefore assume that \(V_1\) and \(V_2\) are distinct.

Use the following fixed connected search. Starting from the exterior,
repeatedly add the lexicographically first unvisited current-lattice
vertex adjacent in the underlying graph to the explored region, and
reveal its entire current cluster. For the wired family start from the
completely exposed wired layer; the same rule interleaves its residual
chambers. The explored union remains connected to the original boundary.
For the two fixed marks count the surrounding odd rims already revealed,
and stop when either counter first reaches \(r\). This rule uses only the
revealed history; it never asks for the identity of a future hole.
If neither counter reaches \(r\), give the stop a cemetery value;
on the event under consideration it is eventually active.
Adjoin the retained stopped states for all these countably many fixed
rules to the joint subsequence. By \cref{prop:exploration-thin}, each
has the stated component laws conditional on its own limiting stopped
state. The future supplier relation is not included in that state.

The two marks lie in one pointed parent at depth \(r-1\). A compact
path inside that parent joining them is eventually inside the same
discrete parent, by pointed convergence. This implication does not
assert the converse through a pinch. The exterior seed must cross
all of their shared enclosing odd rims before it can reach either
\(r\)th rim. Suppose the stop discovers the rim of \(V_1\) first.
The other mark has then seen exactly the common \(r-1\) odd jumps;
any intervening even holes contribute zero. Its stopped value is the
common parent label \(b\).

The newly revealed supplier and its sealing boundary put
\(\partial V_1\) in the limiting explored set \(A\). Hence \(C\)
meets \(A\). If a connected cluster meeting \(A\) entered any component
of the complement, it would also meet that component's boundary.
\Cref{prop:no-crossing}, applied at the appropriate canonical
generation and after subtracting the parent label, forbids this.
Therefore \(C\subset A\), in particular \(\partial V_2\subset A\).

The entire discrete interior of the second odd hole is still
unsearched. A seed vertex could enter it only after meeting its
enclosing supplier, which would have raised the second counter to
\(r\). A different current cluster inside that hole is trapped there
by the supplier and its sealing rim, so it could have been discovered
only from a seed vertex in the same interior. Thus the whole explored
set, not only the seed, misses the second discrete hole. Every compact
subset of \(V_2\), with a slightly larger compact neighborhood,
eventually lies in that untouched hole. It follows that \(A\) misses
\(V_2\). Since \(\partial V_2\subset A\), the domain \(V_2\) is
itself a stopped complementary component.

Its stopped conditional law has flat label \(b\). The labelled odd
carpet gives the same \(V_2\) the label \(b+2a\epsilon(i)\), which
differs from \(b\). Both are conditional descriptions of the same
restricted field on the same pointed domain, so
\cref{lem:recover-flat-label} makes this impossible. The conditional
laws were used at the fixed observable stop, not after conditioning
on the hypothesis of distinct suppliers. This proves the inclusion
at every finite odd depth, including clusters reached after any
number of even restarts.

\begin{figure}[t]
\centering
\begin{tikzpicture}[font=\small,>=Latex]
 \draw[black!45] (0,0) rectangle (6,3.2);
 \node[anchor=north west] at (.1,3.1) {limiting explored set \(A\) (gray)};
 \draw[gray!55,line width=5pt] (0,1.45)--(1.25,1.45);
 \draw[gray!55,line width=5pt] (1.9,1.45) circle (.65);
 \draw[gray!55,line width=5pt] (4.3,1.45) circle (.65);
 \draw[gray!55,line width=5pt] (2.55,1.45)--(3.65,1.45);
 \draw[teal!70!black,line width=1.3pt] (1.9,1.45) circle (.65);
 \draw[teal!70!black,line width=1.3pt] (4.3,1.45) circle (.65);
 \draw[teal!70!black,line width=1.3pt] (2.55,1.45)--(3.65,1.45);
 \draw[->,thick] (.45,1.9)--(.85,1.5);
 \node[anchor=west] at (.15,2.1) {seed};
 \node at (1.9,1.45) {\(V_1\)};
 \node at (1.9,.45) {observed rim};
 \node at (4.3,1.45) {\(V_2\)};
 \node at (4.3,.45) {untouched};
 \node[teal!70!black] at (3.1,2.35) {canonical \(C\subset A\)};
 \node[draw,rounded corners,align=center,text width=4.7cm]
   (same) at (9.1,2.55) {same pointed pair \((V_2,h|_{V_2})\)};
 \node[draw,rounded corners,align=center,text width=2.05cm]
   (stop) at (7.7,1.35) {stopped label\\\(b\)};
 \node[draw,rounded corners,align=center,text width=2.05cm]
   (carpet) at (10.5,1.35) {odd-carpet label\\\(b+2a\epsilon(i)\)};
 \draw[->] (same.south west)--(stop.north);
 \draw[->] (same.south east)--(carpet.north);
 \node[align=center,text width=5cm] at (9.1,.25)
   {circle averages recover one label};
\end{tikzpicture}
\caption{The hypothetical distinct-supplier case at the fixed marked
stop. The first rim puts \(\partial V_1\) in \(A\); no crossing then
puts the whole canonical cluster \(C\), including \(\partial V_2\),
in \(A\), while the discrete interior of \(V_2\) remains untouched.
Its two proposed labels contradict \cref{lem:recover-flat-label}.
The strokes depict inclusion and connectedness, not area or distances.}
\label{fig:visible-partition}
\end{figure}

To exclude extra identifications, use the relative signs. Conditional
on the unsigned lattice data \((N,s)\), the signs of different free
current clusters are independent and fair, and one cluster has the
same sign at all its odd holes. The validity of a marked address is
also unsigned. With our zero sign and false relation at \(\dagger\),
\cref{lem:finite-cuts} therefore gives, for every fixed pair of
addresses in the same current family,
\begin{equation}
 \E[\epsilon_\delta(i)\epsilon_\delta(j)]
       =\P(i\sim_\delta j).
 \label{eq:partition-sign-discrete}
\end{equation}
The canonical fair-sign statement in \cref{lem:canonical-geometry}
gives the same identity with \(\epsilon\) and \(\sim\) for each fixed
pair in that same family. Joint
labelled-carpet convergence passes the bounded sign product to the
limit, while the retained relation bits pass the right side to
\(\P(i\sim_*j)\). Hence
\[
 \P(i\sim_*j)=\E[\epsilon(i)\epsilon(j)]=\P(i\sim j).
\]
The inclusion already proved makes
\(\ind_{\{i\sim_*j\}}-\ind_{\{i\sim j\}}\) nonnegative. Its
expectation is zero, so it vanishes almost surely. This argument
uses global fixed addresses and does not require equality of
discrete parent components to be continuous through a pinch.
The countable intersection over all pairs proves the asserted
partition equality; the relative signs are already the limits
of the marked labels.

Finally choose, for each canonical cluster, one address \(i\)
supplying one of its odd holes. Every other marked odd hole of
that cluster has the same lattice supplier as \(i\) for all
sufficiently small mesh. The pointed-boundary observation above
therefore puts all its odd-hole boundaries in every Hausdorff
limit of \(C_\delta(i)\). Their union is dense in the canonical
cluster by \cref{lem:canonical-geometry}. This proves containment.
For finitely many distinct canonical clusters, the corresponding
relation bits are eventually false, so their suppliers are
eventually distinct.
\end{proof}

Fix an enumeration of the marked addresses. For each canonical cluster,
choose its first supplying address in this enumeration and call the
associated sequence \(C_\delta(i)\) its \emph{visible precursor}.
This is auxiliary information retained for the proof; the matching
topology itself forgets multiplicity.

\begin{corollary}[All-generation local finiteness]
\label{cor:canonical-local-finiteness}
Almost surely, the entire canonical free cluster family has only
finitely many members of diameter greater than any given
$\eps>0$. The same holds for the wired family, including all
descendants. These bounds are compatible with the uniform tight
macroscopic lattice counts.
\end{corollary}

\begin{proof}
Consider any finite family of $M$ canonical clusters of diameter
greater than $\eps$. In each cluster choose two points more than
$\eps$ apart. Density of
odd rims gives two nearly as separated rim points, witnessed by
finitely many nearby odd holes and dense interior test points.
The individual holes need not have diameter comparable to
$\eps$. By \cref{prop:partition}, the two witnesses in each
cluster have one precursor, and the $M$ precursors are distinct.
For sufficiently small mesh each has diameter greater than
$\eps/2$. This finite choice involves only finitely many odd
depths and marks.

Thus, on the joint representation, the event of at least $M$
canonical clusters of diameter greater than $\eps$ implies at
least $M$ lattice clusters of diameter greater than $\eps/2$
eventually. Finite choices from the countable canonical family
and the dense test list suffice; no measurable choice from an
uncountable family is needed. Fatou's lemma and
\cref{cor:global-cluster-counts} make the probability of this event
tend to zero as $M\to\infty$. Countably many rational $\eps$
give the result. The same finite-witness argument can be performed
in the wired descendants, with the boundary cluster counted once.
This proof does not assume nested all-generation local finiteness
in advance.
\end{proof}

\subsection{A flat region cannot contain a revealed crossing}

We give the analytic observation that detects cluster pieces with
no visible jump. It concerns a distributional conditional mean,
not merely a harmonic function on the open complement.

\begin{lemma}[Flat patches]
\label{lem:flat-patch}
Let $h$ be a zero-boundary GFF in a bounded simply connected domain
$O$. Suppose an exploration with information $\cF_A$ describes a
boundary-connected closed set $A$ and a $\cF_A$-measurable distribution $h_A$ such that,
conditionally on $\cF_A$, the residual $h^A=h-h_A$ is the GFF on
$O\setminus A$ extended by zero through $A$, as a distribution on $O$.
The complementary components are assumed simply connected. Extra
randomness may be included in the exploration. For every
deterministic ball $B\Subset O$, almost surely on the event that
$h_A$ is one constant distribution throughout $B$, one has
$A\cap B=\varnothing$. If the conditional decomposition is given
only for tests supported in a neighborhood of $B$ and on a
$\cF_A$-measurable event $\Lambda$, the conclusion holds on $\Lambda$.
\end{lemma}

\begin{proof}
Fix a nonnegative smooth unit-mass mollifier $\rho$ supported in
the unit disk, and put $h_u=h*\rho_u$. For a nonnegative
$b\in C_c^\infty(B)$ set
\begin{equation}
 X_u=u^{-1}\int b(x)\exp(i\beta h_u(x))\,\dd x,
             \qquad \beta^2>2.
 \label{eq:flat-oscillatory}
\end{equation}
On the compact support of $b$, the normalized Green kernel gives
\[
 \E\exp\bigl(i\beta(h_u(x)-h_u(y))\bigr)
 \leq C\left(\frac{u}{|x-y|\vee u}\right)^{\beta^2}.
\]
Consequently
\begin{equation}
 \E|X_u|^2
 \leq C u^{-2}\int_0^{O(1)}
       \left(\frac{u}{r\vee u}\right)^{\beta^2}r\,\dd r
 \leq C.
 \label{eq:oscillatory-L2}
\end{equation}
For a smooth Gaussian exponential cylinder $Y=\exp(ih(\phi))$,
the Gaussian characteristic function and the bounds
$\Var(h_u(x))=\log(1/u)+O(1)$ and
$|\Cov(h_u(x),h(\phi))|\leq C_\phi$ imply
\[
 |\E[X_uY]|\leq C_\phi u^{\beta^2/2-1}\longrightarrow0.
\]
The linear span of these cylinders is dense in $L^2$ of the field.
Together with \eqref{eq:oscillatory-L2}, this proves
$X_u\rightharpoonup0$ weakly in $L^2$. It also proves weak
convergence on the full probability space with extra randomness:
replace any $L^2$ test variable by its conditional expectation
given $h$.

We next bound the conditional \emph{mollified} variance near $A$.
Let $U$ be a simply connected component of $O\setminus A$ and
write $r=|f^{-1}(w)|$, where $f:\mathbb D\to U$ maps zero to $z$.
The conformal Green formula and the growth and quarter estimates
proved or cited in \cref{sec:conformal-growth} give
\[
 G_U(z,w)=-\log r,
 \qquad |z-w|\leq\operatorname{cr}_U(z)\frac{r}{(1-r)^2}.
\]
For $r\geq1/2$ the Green function is at most $\log2$; for
$r<1/2$ the second inequality gives
\[
 G_U(z,w)\leq\log2+
       \log^+\frac{4\operatorname{cr}_U(z)}{|z-w|}.
\]
Koebe's quarter theorem gives
$\operatorname{cr}_U(z)\leq4\dist(z,\partial U)$. If
$\dist(x,A)\leq u$ and $z$ belongs to the mollifier at $x$, then
$\dist(z,\partial U)\leq2u$. Different components have zero
covariance, so the killed Green kernel satisfies on the mollifier
supports
\begin{equation}
 G_{O\setminus A}(z,w)
       \leq\log2+\log^+\frac{32u}{|z-w|}.
 \label{eq:cut-green-bound}
\end{equation}
Rescaling $z=x+u\zeta$, $w=x+u\omega$ and integrating the
locally integrable logarithm against the two mollifiers yields
\begin{equation}
 \Var(h_u^A(x)\mid\cF_A)\leq C_\rho
       \quad\text{if }\dist(x,A)\leq u.
 \label{eq:cut-mollified-variance}
\end{equation}
This is uniform in the exploration. It does not assert a finite
pointwise variance for the GFF itself.

Choose $B_0\Subset B_1\Subset B$ and take $b=1$ on $B_1$.
Let $E$ be the event that $h_A=m$ throughout $B$ and
$A\cap\overline B_0\ne\varnothing$. This is an exploration-
measurable event: obtain $m$ by testing against one fixed smooth
unit-mass function, and check equality of distributions on a
countable dense test family. In the localized version replace $E$
by $E\cap\Lambda$. On $E$, for small $u$,
\begin{equation}
 e^{-i\beta m}\E[X_u\mid\cF_A]
   =u^{-1}\int b(x)
      \exp\left(-\frac{\beta^2}{2}
          \Var(h_u^A(x)\mid\cF_A)\right)\,\dd x.
 \label{eq:flat-positive-integrand}
\end{equation}
All integrands are nonnegative. Since $A$ is connected to
$\partial O$ and meets $B_0$, its $u$-neighborhood in $B_1$ has
area at least $cu$, with deterministic $c>0$. To verify this
without assuming path connectedness, use the distance from a
point in $A\cap\overline B_0$: its image on the connected set
$A\cup\partial O$ is an interval. Choose points at radial spacings
$3u$ up to a fixed positive radius small enough that they lie in $B_1$.
Since $B_1\Subset O$, these points cannot lie on $\partial O$ and
therefore belong to $A$. Their disks of radius $u/2$ are
disjoint, lie in $B_1$, and are within distance $u$ of $A$.
There are at least $c/u$ of them.

Using \eqref{eq:cut-mollified-variance} in
\eqref{eq:flat-positive-integrand} therefore gives
\[
 \E[\ind_E e^{-i\beta m}X_u]\geq c'\P(E).
\]
The left side tends to zero by weak $L^2$ convergence, since
$\ind_Ee^{-i\beta m}$ is a fixed bounded random variable.
Thus $\P(E)=0$. A countable exhaustion of $B$ proves the lemma.
For a random initial component one first conditions on it and
subtracts its initial constant, then uses rational balls with
positive clearance. The conclusion is still a probability-zero
statement, so no bound uniform over all such domains is needed.
\end{proof}

\subsection{Open hole witnesses and the discovered mean}

A cluster with no limiting odd hole might still have many small discrete
holes. The next observation shows that those holes cannot retain positive
area away from the limiting cluster.

\begin{lemma}[Open witnesses for hole contributions]
\label{lem:hole-witnesses}
Let \(O\) be open. Suppose compact sets \(C_n\subset X\) converge in
Hausdorff distance to \(C\), and \(C\cap O\) has zero area. Let
\(U_n\subset O\) be open, with actual relative topological boundary
contained in \(C_n^{\eta_n}\), where \(\eta_n\to0\).
If the indicators of \(U_n\) converge on a dense subset of
\(O\setminus C\) to those of a union \(U\) of components of
\(O\setminus C\), then
\[
 \ind_{U_n}\longrightarrow\ind_U\quad\text{in }L^1_{\rm loc}(O).
\]
In particular, if every nonempty pointed open witness is absent,
the area of \(U_n\) on each compact subset of \(O\) tends to zero.
\end{lemma}

\begin{proof}
For \(x\notin C\), choose a small ball \(B_x\Subset O\setminus C\).
Its closure misses \(C_n^{\eta_n}\) for all large \(n\). Membership
in \(U_n\) is then constant throughout \(B_x\); otherwise the connected
ball would meet its actual relative boundary. A dense test point in
\(B_x\) determines this constant eventually. The indicator of \(U\)
is also constant there, since \(B_x\) lies in one component of
\(O\setminus C\). Thus the indicators converge at every point outside
\(C\). Dominated convergence and the zero area of \(C\) prove the claim.
\end{proof}

We now specify the sweeps used for complete support identification.
They are performed for each current family; the family label is
suppressed in the sweep notation.
For \(e\in\{\pm e_1,\pm e_2\}\), visit current-lattice vertices in
increasing order of \(e\cdot v\), completing every earlier lattice
column before beginning the next one. Within a column use a fixed
lexicographic order. The virtual exterior of the free graph is the
initial connected attachment. In the wired case first reveal the
boundary cluster completely and use it, together with the exterior,
as the attachment for all residual free chambers. When a vertex is
visited, add the edge to its neighbor in the preceding column and
reveal its entire current cluster. That neighbor has already been
visited or belongs to the initial filled region; it is never an
unsearched vertex farther ahead. The seed edges need not be current
edges. Thus the initial region, the added seed, and all revealed
clusters stay connected to the original boundary, and the exact cuts
in \cref{lem:finite-cuts} apply throughout.

We use two countable families of observable stops. A \emph{column stop}
ends after all vertices with \(e\cdot v\leq t\) have been visited, for
a rational \(t\). For rational \(b>0\), a \emph{recorded stop} ends just
after the \(k\)th newly discovered free cluster of diameter at least
\(b\), for a fixed \(k\geq1\). If that discovery never occurs, give
the stop a cemetery value. The rank \(k\) counts only these recorded
discoveries; arbitrarily many smaller clusters may precede them and
remain part of the explored set.

Write \(W_{0,\delta}\) for the initial revealed region: the boundary
perimeter in the free case and the wired current support in the wired
case. Its limit is respectively \(W_0=\partial D\) or
\(W_0=C_\partial\). Write \(S_\delta\) for the added moving seed and
\(A_{\delta,\tau}\) for the full explored set at a fixed stop \(\tau\).
For the wired family this is one composite exploration of the original
field, across all initial chambers. It does not choose a new lattice
exploration from a limiting hole that might have separated from a
shared discrete precursor.

Fix rational compact rectangles
\(Q\subset\operatorname{int}Q^+\Subset D\) and a rational \(r>0\).
On the event
\[
 \dist(Q^+,W_{0,\delta}\cup S_\delta)\geq r,
\]
call a discovered free cluster a \emph{contributor to \(Q^+\)} if
its support meets \(Q^+\) or its union \(\mathcal O_\delta(C)\) of
odd component faces meets \(Q^+\). List these clusters in their
relative discovery order, as
\(C_{\delta,1},\ldots,C_{\delta,L_\delta}\).
A listed cluster contains a seed contact. If its support meets \(Q^+\),
it has diameter at least \(r-O(\delta)\). If an odd hole meets \(Q^+\),
choose a point of that hole there. The cluster surrounds that point
and contains a seed contact at distance at least \(r\).
A point in a bounded complementary component lies in the convex hull
of the surrounding compact set, so the latter's diameter is at least
that distance, again up to the lattice boundary convention.
Therefore, for all sufficiently fine meshes,
\begin{equation}
 L_\delta\leq N_{\delta,0}(r/3),
 \label{eq:contributor-domination}
\end{equation}
where \(N_{\delta,0}(b)=\mathcal N_{\delta,\mathrm{initial}}(b)\)
counts all free clusters of diameter at least \(b\) in the initial
chamber family, as in \cref{eq:spatial-contributor-domination}.
\Cref{cor:global-cluster-counts} gives a uniform tail for this right
side conditional on the initial cut data. This is a pathwise bound
valid at every observable stop; it is not a tail bound after fixing
the completed history at that stop.

For these countably many fixed stops, retain the stopped states of
\cref{prop:exploration-thin}. For each fixed choice of sweep, stop,
\(Q,Q^+\), and \(r\), retain also the active clearance flag and,
when active, the integer \(L_\delta\),
the ordered supports, their relative signs, and the membership bits
of their odd-hole unions at all dense marks. Give the list a
cemetery value when inactive. The integer is tight by
\cref{eq:contributor-domination}, and a fixed-length list has compact
support coordinates in \(\mathscr H\). A further diagonal extraction
therefore makes \(L_\delta\) eventually equal to a finite \(L\) on
the joint representation, with limits \(C_{*,1},\ldots,C_{*,L}\).
The order is reindexed by \(1,\ldots,L_\delta\); raw discovery indices
are not asserted to be tight. Retain also, on the global representation,
the equality bits between each list entry and each marked supplier.
A missing list entry or invalid address has equality bit zero.
A list entry is \emph{visible} if it eventually equals a supplier at
some fixed valid marked address, and \emph{invisible} otherwise.
By \cref{prop:partition}, a visible entry has one associated canonical
cluster, denoted \(C_j^{\rm can}\).

We retain one further kind of auxiliary list on the global
representation. For each fixed \((e,b)\), let \(D_\delta(e,b)\) be
the number of recorded discoveries during the complete sweep, and
write their ordered supports as
\(\widehat C_{\delta,1}^{e,b},\ldots,
\widehat C_{\delta,D_\delta(e,b)}^{e,b}\).
The bound \(D_\delta(e,b)\leq N_{\delta,0}(b)\) makes the integer
length tight. A further diagonal extraction makes it eventually
equal to a finite \(D(e,b)\), and makes each recorded support converge
in \(\mathscr H\); missing ranks have the cemetery value. Retain also
the equality bits between every record and every marked supplier,
set to zero for a missing rank or invalid address.
They are eventually constant on the joint representation.
Call a stabilized record visible if one such bit is eventually one
for a fixed valid address, and invisible otherwise. These lists keep
their order and multiplicity even if several supports have the same
limit. Their lengths are auxiliary coordinates and are not identified
with counts obtained by cutting off the limiting encoded family at
diameter \(b\).

Only coordinates measurable at the fixed stop generate its limiting
sigma-field \(\mathscr E_\tau\): the explored set, seed, observed labels,
and discovered list data. The global representation also carries
future marked suppliers, complete recorded lists, and the canonical
field. They are not added to \(\mathscr E_\tau\) when applying the
stopped conditional laws. In later uses \(r\) is chosen below a
positive limiting clearance, so the discrete clearance flag is
eventually active.

\begin{proposition}[The discovered mean]
\label{prop:discovered-mean}
Let \(\tau\) be one of the fixed stopped sweeps above, with limiting
explored set \(A_\tau\), and write
\(M_\tau=\E[h\mid\mathscr E_\tau]\) as a distribution on tests.
On each connected compact localization \(Q\subset\operatorname{int}Q^+\)
above at positive distance from \(W_0\) and the moving seed,
\(M_\tau\) is the finite signed sum
\begin{equation}
 M_\tau=b_0+2a\sum_{\substack{j\text{ visible}\\
                       \text{in the localized list}}}
       \epsilon_j\ind_{\mathcal O(C_j^{\rm can})}.
 \label{eq:canonical-discovered-mean}
\end{equation}
Here \(b_0=0\) for the original free family. For the wired composite
exploration, \(b_0\in\{-a,a\}\) is the label of the initial pointed
wired hole containing the region. The sum is finite on each such
localization because its contributor list is finite. An invisible
entry has zero contribution. Consequently \(M_\tau\) is a constant
distribution on every ball missing \(W_0\), the moving seed, and the
canonical supports associated with its local contributors.
\end{proposition}

\begin{proof}
Work first with \(Q\subset\operatorname{int}Q^+\) and a positive
clearance as above, on an outcome where the list length has stabilized.
Every \(C_{*,j}\) is contained in \(A_\tau\). By
\cref{prop:exploration-thin}, it has zero area in
\(\operatorname{int}Q^+\), and the actual relative boundary of
\(\mathcal O_\delta(C_{\delta,j})\) lies within \(O(\delta)\) of
\(C_{\delta,j}\).

We identify membership at a dense mark \(q\in\mathcal Q\cap
(\operatorname{int}Q^+\setminus C_{*,j})\). If \(q\) belongs to the
odd-hole union along infinitely many meshes, its positive distance
from \(C_{*,j}\) gives a fixed disk about \(q\) inside that union
on those meshes. The relevant odd depth is bounded by \(L\).
Indeed every enclosing odd ancestor is a distinct cluster: after
entering its odd hole, later clusters lie strictly inside that hole.
The exterior-connected seed must discover each such ancestor before
\(C_{\delta,j}\), and its odd-hole union contains \(q\). It is
therefore another entry of the same list for \(Q^+\).
After passing to a further subsequence, the depth is one fixed
\(r\leq L\), and the representative of \(q\) lies in the fixed disk.
Thus \(C_{\delta,j}\) supplies the marked address
\((\iota,q,r)\). The stabilized equality bit makes the entry visible,
and \cref{prop:partition} identifies its canonical cluster and odd
hole.

Conversely, suppose the entry is visible and \(q\notin C_{*,j}\)
lies in an odd hole of \(C_j^{\rm can}\). That hole has the marked
address \((\iota,q,r)\) at a finite depth. The partition equality
makes its supplier equal to the list entry for all sufficiently
small mesh, so \(q\) eventually belongs to
\(\mathcal O_\delta(C_{\delta,j})\). These two implications show
that the membership indicators converge at every dense mark off
\(C_{*,j}\). In the visible case their limit is
\[
 \ind_{\mathcal O(C_j^{\rm can})\setminus C_{*,j}},
\]
and in the invisible case it is zero. This also rules out
oscillation between membership and nonmembership: any infinite
membership subsequence supplies a fixed visible address, after
which the converse gives eventual membership.

The displayed visible set is a union of components of
\(\operatorname{int}Q^+\setminus C_{*,j}\). In fact
\cref{prop:partition} gives \(C_j^{\rm can}\subset C_{*,j}\), and
the boundary of every canonical odd hole is contained in
\(C_j^{\rm can}\); a component of the smaller complement cannot
cross such a boundary. Removing \(C_{*,j}\) changes no area on this
region. We may therefore apply \cref{lem:hole-witnesses} to each
list entry. Its odd-hole indicators converge in \(L^1(Q)\) to the
canonical odd-hole indicator, or to zero for an invisible entry.
Thus even a collection of long or numerous collapsing holes has
vanishing area unless it supplies an open marked witness.

The signs have stabilized and the list is finite, so the discrete
signed sum in \cref{eq:discrete-discovered-mean} converges in
\(L^1(Q)\) to the sum in the statement. For a visible entry its sign
is the canonical sign because it eventually supplies a fixed marked
address. In the wired case the connected
rectangle \(Q^+\), being disjoint from \(W_0\), lies in one pointed
initial wired hole. Compact containment puts it in one discrete
initial hole for fine mesh, even if that hole pinches elsewhere.
Its initial label is therefore one constant \(b_0\in\{-a,a\}\)
on \(Q^+\). This supplies the additional baseline term.
Remove a deterministic truncation \(L\leq M\) using
\cref{eq:contributor-domination} and the initial count tail.
The uniform higher logarithmic moments in \cref{prop:log-moments}
give uniform integrability when passing tested conditional means.

Finally, \cref{prop:exploration-thin} identifies the full conditional
mean on these tests, including across \(A_\tau\), by its neighborhood
removal argument. Hence the ordinary function just obtained is
\(M_\tau\), with no extra distribution on the explored set.
On a connected ball avoiding \(C_j^{\rm can}\), its odd-hole indicator
is constant. The baseline is also constant on a ball avoiding \(W_0\).
The localized list contains only finitely many entries, regardless
of the diameters of their associated canonical clusters. This proves
the last assertion.
\end{proof}

\subsection{Complete matching, including invisible clusters}

\begin{proposition}[Matching of all positive-diameter clusters]
\label{prop:complete-matching}
Every positive-diameter Hausdorff limit of lattice free-current clusters
is one retained canonical cluster, and every retained canonical cluster
is the limit of its visible precursor. No distinct invisible precursor
has a positive-diameter limit. The same assertions hold for the free
primal descendants after the wired layer, jointly with that layer and
the height. Consequently
\[
 F_*^{\mathrm f}=F(\cC^{\rm f}),\qquad
 F_*^{\mathrm w}=F(\cC^{\rm w}),
\]
and both collections converge in the matching topology of
\cref{thm:main}.
\end{proposition}

\begin{proof}
\emph{Every point is on a canonical support or the initial wall.}
Let \(C_\delta\) be lattice free clusters with a positive-diameter
Hausdorff limit \(C_*\). The limit is compact and connected.
Take \(x\in C_*\setminus W_0\) and a distinct point \(y\in C_*\).
Choose an axial direction \(e\) and a rational \(t\) with
\(e\cdot y<t<e\cdot x\). At the fixed column stop \(\tau(e,t)\),
the cluster \(C_\delta\) has been discovered for all sufficiently
small mesh, because it has a point near \(y\) in an earlier column.
A small rational ball about \(x\) lies ahead of the seed and has
positive clearance from \(W_0\). The same clearances hold at fine mesh.
Choose \(Q\subset\operatorname{int}Q^+\) around that ball and a rational
lower clearance \(r\). The target cluster is an entry in the finite
localized list at this already specified stop.

Let \(\mathcal U_\tau\) be the finite union of the canonical supports
associated with that list's visible entries. If \(x\notin\mathcal U_\tau\),
a smaller rational ball \(B\) around \(x\) misses \(\mathcal U_\tau\),
\(W_0\), and the moving seed. By \cref{prop:discovered-mean},
\(M_\tau\) is a constant distribution on \(B\).
Let \(\Lambda_{Q^+}=\{\dist(Q^+,W_0\cup S)>0\}\), the
stopped-measurable clearance event for this localization. In the
present initial domains it is the event in
\cref{eq:stopped-full-characteristic}. That identity gives, for every
real smooth test \(\varphi\) supported in \(Q^+\),
\begin{equation}
 \ind_{\Lambda_{Q^+}}
 \E\!\left[e^{i(h-M_\tau)(\varphi)}\mid\mathscr E_\tau\right]
 =\ind_{\Lambda_{Q^+}}\exp\!\left\{-\tfrac12
       G_{D\setminus A_\tau}(\varphi,\varphi)\right\}.
 \label{eq:full-residual-use}
\end{equation}
It holds for all linear combinations of finitely many tests, with the
killed GFF extended by zero across \(A_\tau\). It therefore supplies
the local residual law required by \cref{lem:flat-patch}, including
the mollifiers whose supports cross the explored set. Since \(A_\tau\) is
boundary-connected and the complementary components are simply connected, that lemma
says \(A_\tau\cap B=\varnothing\). But \(C_\delta\) was revealed at
this stop, so \(C_*\subset A_\tau\), a contradiction.

The choice of \(B\) may depend on the later canonical data. Its
existence implies one of the countably many
\(\mathscr E_\tau\)-measurable zero events that \(M_\tau\) is constant
on a fixed rational ball and \(A_\tau\) meets it. The conditional law
used in each of these events is the one at the fixed stopped state.
All directions, column thresholds, balls, and localizations come from
countable families, and each localized list has finitely many entries.
The preceding conclusion therefore holds simultaneously for every
possible subsequential \(C_*\) and every \(x\in C_*\setminus W_0\).
We have proved
\begin{equation}
 C_*\subset W_0\ \cup\!\bigcup_{C\in\cC_{\rm can}}C,
 \label{eq:countable-cover}
\end{equation}
where \(\cC_{\rm can}\) is the complete canonical free family under
consideration. For wired descendants the argument used the single
composite exploration of the original \(h\). A ball away from \(W_0\)
lies in one pointed wired hole and has a constant initial label, even
when its discrete initial chamber pinches into several limiting holes.

The sets in the cover are compact and pairwise disjoint, including
\(W_0\): this follows from \cref{lem:canonical-geometry} and the fact
that each retained cluster lies strictly inside its parent. They form
a countable family. Their intersections with \(C_*\) are therefore
a disjoint closed cover of the compact connected set \(C_*\).
Sierpi\'nski's continuum partition theorem \cite{Sierpinski1918},
with a complete proof in \cite[Chapter~V, Theorem~5.10]{Freiwald2014},
says that at most one member of such a cover is nonempty.
Consequently \(C_*\) is contained in one canonical cluster or entirely
in \(W_0\).

We exclude the latter case using the actual chamber wall.
Let \(\Gamma_\delta\) be the component-side open \(W\) perimeter of
the initial chamber containing \(C_\delta\), with the medial resolution
at pinches. In the wired case \(W_{0,\delta}\) is a current support,
whereas the spatial estimate uses the attached matching wall.
\Cref{lem:chamber-rim-distance} gives
\[
 \sup_{z\in C_\delta}\dist(z,\Gamma_\delta)
 \leq \sup_{z\in C_\delta}\dist(z,W_{0,\delta})+O(\delta).
\]
Its first-exit proof uses the resolved chamber face: a segment from
an adjacent chamber point toward the revealed support first meets
this chamber's own medial rim, including at a pinch. The free case
uses the polygonal boundary.
If \(C_*\subset W_0\), Hausdorff convergence of \(C_\delta\) and
\(W_{0,\delta}\) makes the right side tend to zero.

Choose one fixed buffered square \(Q_{\rm all}\) containing \(X\)
inside a larger regular ambient box,
and use the event \(\Xi_\delta(\tau_0;b,\eta)\) of
\cref{cor:indexed-no-shadowing} in this square, at the initial stop
\(\tau_0\). Its connected cell region contains the initial revealed
support up to \(O(\delta)\). If the limit above has diameter greater
than \(2b\), then for every fixed \(\eta>0\), at all sufficiently fine
meshes \(C_\delta\) has a connected passage of diameter at least \(b\)
in \(Q_{\rm all}\) inside that \(\eta\)-collar. It is still unsearched
at \(\tau_0\). Thus the existence of such a limit forces membership in
\(\liminf_n\Xi_{\delta_n}(\tau_0;b,\eta)\) for every fixed \(\eta\).
For any deterministic \(\eta_m\downarrow0\), Fatou's inequality and
the bound in \cref{eq:spatial-frozen-shadow}, integrated over the
initial cut data, give
\[
 \P\!\left[\bigcap_m\liminf_n
       \Xi_{\delta_n}(\tau_0;b,\eta_m)\right]
 \leq \inf_m\limsup_{\delta\downarrow0}
       \omega_{Q_{\rm all},b}(\eta_m,\delta)=0.
\]
The collar width is fixed before each mesh limit; its possible
dependence on the future Hausdorff limit is used only for the
eventual inclusion. Countably many rational \(b>0\) exclude every
positive-diameter limit contained in \(W_0\). We conclude that
\(C_*\) is contained in one canonical cluster.

\emph{Visible precursors.}
For the visible precursor of a canonical cluster \(C\),
\cref{prop:partition} puts \(C\) in every Hausdorff limit.
The preceding containment puts that limit in a single canonical
cluster. Pairwise disjointness forces it to be \(C\) itself.
Compactness of \(\mathscr H\) then shows that the entire visible
precursor sequence converges to \(C\). In particular every canonical
cluster is obtained.

\emph{Invisible precursors.}
Return to any lattice sequence with positive-diameter limit
\(C_*\subset C\). Choose an axial direction \(e\) in which
\[
 w=\max_{z\in C_*}e\cdot z-\min_{z\in C_*}e\cdot z>0,
\]
and fix a rational \(b\in(0,w/4)\). This sequence is eventually among
the \((e,b)\) records. Their length has stabilized to a finite number,
so one fixed rank \(k\) supplies it along infinitely many meshes.
The full sequence \(\widehat C_{\delta,k}^{e,b}\) already has a
Hausdorff limit, which must equal \(C_*\) on that subsequence. Replace
the possibly alternating original sequence by this fixed recorded
sequence. This replacement classifies its support limit; it makes no
assertion that visibility transfers from the original choices.

If this record is visible, \cref{prop:partition} puts its associated
canonical cluster inside \(C_*\subset C\). Pairwise disjointness makes
that cluster \(C\), and the stabilized supplier bits identify the
record with the visible precursor of \(C\) for all fine meshes.
It was already treated. Suppose therefore that the record is invisible.
Write \(C_\delta^{\rm vis}\) for the visible precursor of \(C\).
It too is eventually a record because it converges to \(C\) and
\(\diam C\ge w>4b\). Among the finitely many stabilized ranks, exactly
one has its supplier bit equal to one for the fixed chosen address of
\(C\). Thus \(C_\delta^{\rm vis}\)
has one fixed rank \(j\), and \(j\ne k\). The two ranks give one
eventual discovery order. These complete-list coordinates remain
outside the conditioning sigma-field of every fixed recorded stop.
Smaller discoveries remain in each explored set; their local
contributions are controlled by \cref{eq:contributor-domination}.

Suppose first that the visible precursor is discovered earlier, at
the fixed recorded stop \(\tau\). Its limit is \(C\).
For each fixed \(\eta>0\), the convergences
\(\widehat C_{\delta,k}^{e,b}\to C_*\subset C\) and
\(C_\delta^{\rm vis}\to C\)
imply that eventually the later cluster lies inside the
\(\eta\)-neighborhood of the frozen cell region at \(\tau\), using
the \(O(\delta)\) displacement from \(A_{\delta,\tau}\). It is still unsearched
at \(\tau\) and has a passage of diameter at least \(b\) inside
\(Q_{\rm all}\). Let
\[
 \mathcal X_\delta(b,\eta)=
 \{\exists\,0\le i\le D_\delta(e,b):
              \Xi_\delta(\tau_i;b,\eta)\}
\]
for the fixed \((e,b)\) sweep in that square. The event under
consideration is contained in \(\liminf_n\mathcal X_{\delta_n}(b,\eta)\)
for every fixed \(\eta\). For every integer \(K\),
\cref{eq:spatial-indexed-shadow} and Fatou's inequality yield
\[
 \P\!\left[\bigcap_m\liminf_n
             \mathcal X_{\delta_n}(b,\eta_m)\right]
 \le Ce^{-cK}+(K+1)\inf_m\limsup_{\delta\downarrow0}
          \omega_{Q_{\rm all},b}(\eta_m,\delta)
 =Ce^{-cK}.
\]
Letting \(K\to\infty\) gives probability zero. The spatial bound is
applied to the existential future passage under its allowed history
\(\mathcal H_\tau\), for finitely many fixed indices before the count tail is
removed. Later visibility and limiting containment only give the
eventual inclusion in the fixed-\(\eta\) events; they are not added
to that conditioning.

Suppose instead that the invisible cluster is discovered first, and
stop immediately after that recorded discovery. Let \(t_\delta\) be
the projection of its discovery column. Every earlier column has
already been visited, so the first occupied column of this cluster
satisfies
\[
 \left|t_\delta-
       \min_{z\in\widehat C_{\delta,k}^{e,b}}e\cdot z\right|=O(\delta).
\]
Hausdorff convergence preserves the minimum and maximum of the
projection. Choose a point \(x\in C_*\) whose projection is more
than \(w/2\) above the minimum. A rational ball about \(x\) lies
a fixed positive distance ahead of the moving seed for all fine
meshes. It also has positive clearance from \(W_0\), since
\(C_*\subset C\) and \(C\) is compactly inside its initial parent.
The stopped explored set meets that ball.

Use its localized contributor list at this fixed recorded stop.
The newly discovered record meets the ball, so it is the last entry
of that list and has the fixed index \(L\) after the length stabilizes.
That entry is invisible: a stabilized marked supplier bit equal to
one would also make the record visible. It has zero mean contribution
by \cref{prop:discovered-mean}. Every earlier relevant contributor,
whether its support meets the ball or its odd hole surrounds a point
there, is included in the same finite list. Its associated canonical
support cannot be \(C\): by \cref{prop:partition} that would make
it the already chosen visible precursor of \(C\), contrary to the
discovery order. The earlier visible supports are therefore a finite
family of compact sets disjoint from \(C\). Shrink to a rational ball
around \(x\) missing this union. Its baseline and all remaining
odd-hole indicators are constant, so \(M_\tau\) is constant there.
The full residual law \cref{eq:full-residual-use} applies at this
stop, and \cref{lem:flat-patch} contradicts the hit of the explored
set. As before, the later choice of a rational ball only implies
a stopped-data zero event and adds no conditioning.
This excludes the second order and hence every distinct invisible
record with a positive-diameter limit. Since every positive limit
of an arbitrary lattice sequence was also the limit of a recurrent
fixed record, the classification is exhaustive. There is only one
visible rank for a given canonical cluster. A sequence distinct from
that precursor along infinitely many meshes has a recurrent rank
\(k\ne j\) on those meshes, by finiteness of the record list. If that
rank were visible, its limiting support inside \(C\) would make it
the same chosen supplier and hence rank \(j\) eventually. Thus it is
an invisible recurrent rank, which has just been excluded.

\emph{The collection topology.}
For clarity, outer Hausdorff convergence has the following direct
link with the individual limits just classified. If \(C\in F_*^\iota\)
has positive diameter, then
\(\dist_{\mathscr H}(C,Z)>0\). A nearest member of
\(F(\cC_\delta^\iota)\) approaching \(C\) therefore cannot be a
singleton for fine mesh; it is an actual lattice cluster.
Conversely, every positive-diameter Hausdorff limit of lattice
members belongs to \(F_*^\iota\), because these encoded families
converge and \(F_*^\iota\) is closed. The fixed sweeps and their
finite localized lists made the preceding classification simultaneous,
so it applies to these members even when they are selected from the
outer limit.

It follows that the positive members of \(F_*^{\mathrm f}\) are
exactly \(\cC^{\rm f}\). For \(F_*^{\mathrm w}\), a selected lattice
member has a subsequence of one recurrent type: either the distinguished boundary cluster, whose limit is
\(W_0=C_\partial\), or a free descendant, which has just been classified.
The visible precursors give the reverse inclusion in both cases.
Thus the two encoded limits equal those in the statement.
\Cref{lem:matching-topology} converts this outer Hausdorff convergence
to the stated matching topology. The exclusion of duplicate invisible
precursors is a stronger conclusion retained in the auxiliary lists;
the set-valued metric itself forgets multiplicity.
\end{proof}

\subsection{Completion of the joint limit}

\begin{proof}[Proof of \cref{thm:main}]
By \cref{thm:flat-gff}, $K_\delta$ converges to $h/a$ in every
local Sobolev space $H^s_{\rm loc}$, $s<-1$. The arguments above
can be performed on any simultaneous subsequence of the field,
the wired current, and the free dual current. By
\cref{prop:stopped-tvs}, the distinguished wired boundary
cluster converges in ordinary Hausdorff distance to
$\TVS{a}{a}(h)$. Its remaining pointed domains have the
conditional zero-boundary GFF remainders of that same field.

Apply \cref{prop:complete-matching} to the free dual family in
the original domain and to the free primal descendants in the
composite wired exploration. The positive outer $\CLE_4$ label is $2\lambda$,
and its retained split has absolute values $0,2a$; the negative
one is its reflection. The identified recursion is therefore
exactly \cref{def:free-procedure}, with its sign-sensitive
choice of split, and with every generation retained. Both
families have been identified as canonical functions of the
same limiting $h$.

For completeness, these are measurable collection-valued functions.
Enumerate the recursive clusters by generation and by the first mark
of \(\mathcal Q\) in each generating hole, using a cemetery value for
missing entries. The cluster coordinates in this countable enumeration
are measurable by \cref{lem:canonical-geometry}.
Append \(C_\partial\) as one additional slot for the wired family.
\Cref{cor:canonical-local-finiteness} makes their singleton-augmented
families closed. For every open set \(\mathcal V\subset\mathscr H\),
the event that \(F(\cC^\iota)\) meets \(\mathcal V\) is the union of
the corresponding countably many cluster hit events, together with
the deterministic alternative \(Z\cap\mathcal V\ne\varnothing\).
These hyperspace hit events prove measurability of both encoded factors.
The boundary marker \(\TVS{a}{a}(h)\) is measurable in the same field.

The height and collection tightness give simultaneous joint
subsequences, and every one has the pushforward of the GFF law by
this single measurable construction.
Hence the full mesh limit exists in the product topology,
including the distinguished boundary set. Finally
$H_\delta=K_\delta/2$ and $2a=\pi\sqrt g$, giving precisely
$h/(\pi\sqrt g)$. The parameter inequalities used throughout
include $a=2\lambda$, so the Potts endpoint is part of the
conclusion.
\end{proof}
